
Formally, for constants $k\ge1$ and $\rho\in(0.8,0.9)$, let $r := 2k+1$, 
$L := \lfloor\log_2 n\rfloor$, 
$\mathcal{P} := \{\,x : \mathrm{LO}(x)\ge 0.9n\,\}$, 
$\mathcal{B}_k := \{\,0^{r} 1^{(0.9n-r)} *\,\}$, 
and $\mathcal{D}_\rho := \{\,0^{L} * : |x|_1 \ge \rho n\,\}$.
The function $\mathrm{LTJ}_{k,\rho} : \{0,1\}^n \to \mathbb{R}$ is
\[
\mathrm{LTJ}_{k,\rho}(x)
=
\begin{cases}
	0, & x \in \mathcal{P},\\[2pt]
	\mathrm{LO}\bigl(x_{i>r}\bigr)+r, & x \in \mathcal{B}_k,\\[2pt]
	\textcolor{black}{n+1}, & x \in \mathcal{D}_\rho,\\[2pt]
	\mathrm{LO}(x), & \text{otherwise.}
\end{cases}
\]
\textcolor{blue}{Here $x_{i>r}$ denotes the suffix of $x$ at positions $i > r$, i.e.,
$x_{i>r} := (x_{r+1}, x_{r+2}, \dots, x_n)$.}
 \textcolor{blue}{Note that  $r := 2k+1$ is an odd number, as required for the block-jump prefix $0^r1^{(0.9n-r)}*$
to force an even number of bit-flips to reach $\mathcal{B}_k$, which is
the key parameter controlling the per-iteration success probability
$\Theta(n^{-(r+1)}) = \Theta(n^{-(2k+2)})$. We introduce $r$ as shorthand
to avoid writing $2k+1$ in every exponent throughout the proof.}
\begin{figure}[t]
\centering
\includegraphics[width=0.48\textwidth]{ltj_diagram}
\caption{\textcolor{blue}{The \textsc{LeadingTrapJump}$_{k,\rho}$ landscape (schematic,
$n=30$, $k=1$, $\rho=0.85$).
The x-axis is the \textsc{LeadingOnes} level of the current solution.
From the junction~$J$ (end of the main slope), large mutations can reach
either the block-jump slope~$\mathcal{B}_k$ (leading to the local optimum)
or the deep-prefix plateau~$\mathcal{D}_\rho$ (global optimum).
Both off-slope structures have $\mathrm{LO}(x)=0$; their positions on the
y-axis reflect their fitness values, not their \textsc{LeadingOnes} level.}}
\label{fig:ltj}
\end{figure}

The landscape consists of a long LeadingOnes slope rising from fitness 
$\mathrm{LO}(x)$ up to level $0.9n-1$, followed by a large penalised region.  
Two off-slope structures contain the \textcolor{black}{two sets of optima} but are inaccessible to 
standard hill-climbers.  
The set $\mathcal{B}_k$ begins with the prefix $0^{r}1^{(0.9n-r)} *$ and forms a 
short slope, \new{called block-jump slope,} that \textcolor{black}{deterministically increases toward its unique local optimum 
\textcolor{blue}{$0^{r}1^{(0.9n-r)}$}; reaching 
this prefix is therefore equivalent to entering the leading ones path to the local optimum\textcolor{blue}{
$0^{r}1^{(0.9n-r)}$}.}  
In contrast, the set $\mathcal{D}_\rho$ consists of solutions with prefix 
$0^{L} *$ \textcolor{black}{whose suffix has ones-density at least $\rho$} i.e., $\rho n$ ones, and forms a wide 
plateau of globally optimal fitness, \new{which we call  deep-prefix plateau}.  
\textcolor{blue}{Both structures are hidden behind the penalised barrier (illustrated in
Figure~\ref{fig:ltj}), and the only reliable way to reach them is via large
hypermutation events.}  
The junction $J := \{x : \mathrm{LO}(x)=0.9n-1\}$ acts as a central gateway: 
each ageing epoch that restarts the search will, with high probability, ascend 
the LeadingOnes slope back to $J$ before any off-slope structure is encountered.  
From $J$, the \oneoneOPTIA~using \IPHfcm{}~with \linHD~can perform jumps of the 
appropriate magnitude to reach either the entrance of the block-jump slope $\mathcal{B}_k$ or 
the plateau $\mathcal{D}_\rho$ \rev{i.e., via ageing restarts it will eventually find both optima efficiently.} 
\textcolor{black}{On the other hand, algorithms with lower mutation rates, such as those using SBM will with high probability not be able to flip the logarithmic number of bits required to reach the global optimum, thus have a runtime of $n^{\Omega(\log n)}$ with high probability.}

\begin{theorem} \label{thm:main-tau-final}
The  $(1+1)$~Opt-IA$_{\textsc{FirstBest}}$~with ageing threshold $\tau=\Omega(n^{1+\epsilon})$ using  \linHD~with $c<5/8$  identifies both local and global optima of  $\mathrm{LTJ}_{k,\rho}$ in
	$O(n\tau + n^{2k+3} + n^{\beta+1})$ fitness function evaluations	where $\beta := \log_2 \frac{20}{c}$.
\end{theorem}

To establish Theorem~\ref{thm:main-tau-final} we rely on three ingredients.
First, every ageing epoch returns to the junction $J$ with high probability.
Second, an epoch beginning at $J$ performs sufficiently long jumps to reach the beginning of the block-jump slope with probability $\Theta(n^{-(2k+2)})$ per iteration.
Third, an epoch beginning at $J$ performs a logarithmic jump onto the deep-prefix plateau with probability $\Theta(n^{-\beta})$ per iteration.

Since each event corresponds to reaching either the start of a short slope leading deterministically to its optimum (for $\mathcal{B}_k$) or directly a plateau of global optima (for $\mathcal{D}_\rho$), these per-epoch probabilities govern the entire runtime.

We first need to bound from below the probability that the algorithm visits the junction point after restarts in order to later argue that from the junction point both sets of optima will be accessed.

\begin{lemma} \label{lem:junction-first-density-Bk}
	Starting from any random individual, the $(1+1)$~Opt-IA$_{\textsc{FirstBest}}$~using \IPHfcm{}~with \linHD~\textcolor{black}{and $c<5/8$}  reaches the junction $J = \{x : \mathrm{LO}(x)=0.9n-1\}$ before accepting any individual in $\mathcal{B}_k$ or $\mathcal{D}_\rho$ with high probability.
\end{lemma}

Within $n^2$ iterations, the algorithm performs at least $0.9n$ LeadingOnes 
improvements with high probability.  
Before reaching $J$, deep-prefix entrances require flipping the $L$-bit prefix 
to $0^L$, an event whose probability per iteration is at most $(M(x)/n)^L$.  
Since under \linHD\ the mutation potential satisfies $M(x) \le q^\star n$ for 
some constant $q^\star < 1$ (determined solely by $c$ and $\rho$), this becomes 
$(q^\star)^L = n^{-\alpha^\star}$ for $\alpha^\star := \log_2(1/q^\star) > 2$, 
yielding a total probability of $o(1)$ over $n^2$ iterations.
Thus the slope is climbed to $J$ with high probability.


\begin{proof}
	Let $T:=n^2$. \textcolor{purple}{We will first show that in $T$ iterations with overwhelming probability we sample a solution from  $J$ before we sample one from $\mathcal{B}_k$ or $\mathcal{D}_\rho$.}
%	that unless we sample a solution from $\mathcal{B}_k$ or $\mathcal{D}_\rho$, 
%	we will sample a solution from $J$ in $T$ iterations with overwhelming probability.}
	
	Before $J$, $\mathcal{B}_k$ or $\mathcal{D}_\rho$ is sampled from, we flip the leftmost zero-bit with probability $1/n$ and accept
	immediately by FCM($\ge$).  If $X\sim\mathrm{Bin}(T,1/n)$ counts such improvements then
	$\mathbb{E}[X]=n$ and a Chernoff bound gives
	$
	\Pr\bigl(X<(9/10)n\bigr)\ \le\ e^{-a n}
	$
	for some constant $a>0$. Hence $J$ is reached within $T$ iterations with probability $1-e^{-a n}$ if we condition on that neither $\mathcal{B}_k$ nor $\mathcal{D}_\rho$ is sampled first.
	
	\textcolor{black}{
		Next, we will show that the probability of sampling a $\mathcal{B}_k$ or a $\mathcal{D}_\rho$ solution in $T$ iterations is $o(1)$, starting our argument with $\mathcal{D}_\rho$. 
		}
	
	Let $\alpha:=|x|_1/n$ and $\alpha^\top:=|\text{best}|_1/n$. Since the probability of improvement is at least $1/n$ in the initial leading-ones slope the ageing does not trigger with overwhelming probability unless the current solution is in $J$, $\mathcal{D}_\rho$ or $\mathcal{B}_k$, all of which lead to $\alpha^\top\ge\rho$ after the restart. We make a case distinction according to $\alpha\ge \rho-\xi$ for some constant $\xi\in(0,\rho-\tfrac12]$  such that	$q^\star \;:=\; 2c\,\bigl(1-\rho+\xi\bigr) \;<\; 1/4$.
\textcolor{blue}{This is achievable: choosing $\xi$ sufficiently close to $0$ gives
$q^\star \approx 2c(1-\rho)$; since $\rho > 4/5$ we have $1-\rho < 1/5$,
so $2c(1-\rho) < 2 \cdot \tfrac{5}{8} \cdot \tfrac{1}{5} = \tfrac{1}{4}$
(using $c < 5/8$).  Any fixed $\xi \in (0, \rho - 1/2]$ small enough to
keep $2c(1-\rho+\xi) < 1/4$ then satisfies all requirements.}
	
	\emph{Case 1: $\alpha\ge \rho-\xi$.}
	Using $|x\wedge best|_1\ge |x|_1+|best|_1-n$,
	\[
	\mathrm{HD}(x,best)\ \le\ (2-\alpha-\alpha^\top)\,n\ \le\ 2(1-\rho+\xi)\,n,
	\]
	\[
	M(x)\ \le\ q^\star n\ <n/4
	.
	\]
	$\mathcal{D}_\rho$ definition requires all $L$ leading bits among the $M$ flips; hence
	\[
	\Pr(\text{$\mathcal{D}_\rho$ in this iteration})
	\ \le\
	\Bigl(\tfrac{M}{n}\Bigr)^{L}
	\ \le\
	(q^\star)^{L}
	\ =\
	n^{-\log_{2}(1/q^\star)}.
	\]
	Because $q^\star<1/4$, we have $\log_{2}(1/q^\star)>2$, so $T\cdot n^{-\log_{2}(1/q^\star)}=o(1)$.
	
	\emph{Case 2 : $\alpha\le \rho-\xi$.}
	Since all $\mathcal{D}_\rho$ solutions have $|x|_1\ge \rho n$, at least $\xi n$ net $0\!\to\!1$ flips ($Z-O$) are required within at most
	$cn$ bit-flips. With initial zero/one ratio giving negative drift,
	a Serfling (hypergeometric) tail bound yields
	\[
	\Pr\bigl(Z-O\ge \xi n\bigr)\ \le\ e^{-\gamma n}
	\]
	for some constant $\gamma>0$. Including the necessity to flip all $L$ prefix bits only decreases the probability.
	Summing over at most $T$ iterations remains $e^{-\gamma n}$. Therefore $\Pr(\text{$\mathcal{D}_\rho$ before }J)=o(1)$.
	
	Similarly we will next establish that sampling from $\mathcal{B}_k$ in $T$ iterations  has probability $o(1)$.	Let $r:=2k{+}1$, $S:=\{r{+}1,\dots,\tfrac{9}{10}n\}$, and fix any constant
	$\delta\in(0,\tfrac{9}{10})$.  Consider any iteration with parent
	$x\notin\mathcal{B}_k$ and write $\ell:=\mathrm{LO}(x)$ and
	$v:=|\{i\in S:x_i=0\}|$.  We distinguish two cases depending on $\ell$.
	
	\emph{Case 1: $\ell < (\tfrac{9}{10}-\delta)n$.}
	Then the unconstrained segment
	$S\cap\{\ell{+}1,\dots,\tfrac{9}{10}n\}$ has size at least $\delta n$.
	Conditioned on the history up to the current iteration, the bits in this segment
	are unbiased and independent, hence by Chernoff $v=\Omega(n)$ with probability
	$1-e^{-\Theta(n)}$.  On this event, to enter
	$\mathcal{B}_k=\{0^{r}1^{(0.9n-r)}*\}$ and be accepted under FCM($\ge$), the
	offspring at acceptance must flip all $r$ positions $\{1,\dots,r\}$ and must flip
	all $v=\Omega(n)$ zeros in $S$ to $1$, i.e., it must flip $\Omega(n)$ specified
	positions.  Therefore, regardless of the mutation potential, the per-iteration
	probability of entering $\mathcal{B}_k$ is at most
	$\binom{n}{\Omega(n)}^{-1}=e^{-\Theta(n)}$.  Union-bounding over at most
	$T\le n^2$ iterations yields a total contribution $e^{-\Theta(n)}$.
	The exceptional event $v=o(n)$ itself has probability $e^{-\Theta(n)}$ per
	iteration and is negligible under the same union bound.
	
	\emph{Case 2: $\ell \ge (\tfrac{9}{10}-\delta)n$.}
	Then the set
	$F:=\{r{+}1,\dots,\ell\}$ of leading-one positions that must not be flipped before
	acceptance has size $|F|=\ell-r=\Omega(n)$.  Since $x\notin\mathcal{B}_k$, at
	least one bit in $S$ is $0$, so entering $\mathcal{B}_k$ and being accepted
	requires flipping the $r$ positions $\{1,\dots,r\}$ and at least one further
	specified position in $S$ before any position in $F$ is flipped (otherwise
	$\mathrm{LO}$ decreases and FCM($\ge$) rejects).  Let $R$ denote any fixed choice
	of such $r{+}1$ required positions.  A necessary condition for acceptance into
	$\mathcal{B}_k$ is that all positions in $R$ appear before all positions in $F$ in
	the random order in which the hypermutation flips bit positions.  Since the
	relative order of the $|F|+|R|$ positions in $F\cup R$ is uniform, we obtain
	\[
	\Pr\bigl(\mathcal{B}_k \text{ in this iteration}\bigr)
	\le
	\binom{|F|+r+1}{r+1}^{-1}
%	\le
%	\Bigl(\tfrac{r+1}{|F|}\Bigr)^{r+1}
	=
	O\!\bigl(n^{-(r+1)}\bigr).
	\]
	Union-bounding over at most $T\le n^2$ iterations yields a total contribution
	$O(n^{-(r-1)})=O(n^{-2k})=o(1)$.
	
\end{proof}

\textcolor{purple}{The following two lemmata will condition throughout on the event that during an ageing epoch the junction $J$ is reached.
This event happens with high probability by Lemma~\ref{lem:junction-first-density-Bk}.}

\begin{lemma} \label{lem:BJO-epoch}
	\textcolor{blue}{Let an ageing epoch reach the junction $J$.} %(which occurs 
	%with high probability.
	%by Lemma~\ref{lem:junction-first-density-Bk}; we condition on this event
	%throughout).}
	Then, with probability $\new{p_B}=\Omega(\min\{1,\,\tau\,n^{-(2k+2)}\})$, the epoch performs a jump placing the individual at the beginning of the block-jump slope.
\end{lemma}

The suffix of the individual at $J$ is essentially random relative to the best individual stored from the 
previous epoch. 
With high probability their Hamming distance is at least $n/20$, guaranteeing 
\textcolor{black}{$M(x)= cn/20$} throughout the epoch.  
Under this event, the probability that the first $r+1$ flips of the hypermutation hit exactly the positions defining the prefix $0^{r}1^{(0.9n-r)} *$ is $\Theta(n^{-(r+1)})$.
Evaluating this over $\tau$ iterations yields the stated bound.


\begin{proof}
	At the beginning of the epoch %start %the reference best is the best individual from the
	%previous epoch.  
	the length-\(0.1n+1\) suffix of the new parent \(x\) is still uniformly random,
	so with probability at least \(\tfrac12\)
	\[
	\mathrm{HD}(x,best)\;\ge\;\frac{n}{20}.
	\]
	Conditioning on this event we lose only a factor \(2\).
	Then the mutation potential in every subsequent
	iteration of the epoch satisfies
	\(
	M\ge \tfrac{cn}{20}  \geq r+1.
	\) for large enough $n$.
	
	%\paragraph*{Per-iteration success probability (revised, $M\ge r{+}1$).}
	%Conditioned on $\mathrm{HD}(x,x^\top)\ge n/20$ we certainly have
	%$M=\lfloor c\,\mathrm{HD}\rfloor\ge r+1$.
	A solution in $\mathcal{B}_k$ is accepted as soon as the \emph{first} $r{+}1$ distinct flips
	include \emph{exactly} the set
	\(\{1,\dots,r,\;0.9n\}\)
	in any order; if this occurs, FCM stops immediately and no other
	leading-one bit can be destroyed.
	%
	Since bits are sampled uniformly without replacement, the probability of
	this event is
	\[
	p_{\text{iter}}
	=\prod_{i=0}^{r}\frac{r+1-i}{n-i}
	=\frac{(r+1)!}{(n)_{\,r+1}}
	=\Theta\!\bigl(n^{-(r+1)}\bigr)
	=\Theta\!\bigl(n^{-(2k+2)}\bigr),
	\]
	independent of the exact value of~$M$ (we only need $M\ge r{+}1$ to allow that
	many flips).
	
	%\medskip\noindent
	%\emph{Per-epoch success probability.}
	Since an epoch lasts for \(\tau\) independent iterations, a union bound yields
	\[
	p_B
	\;\ge\;
	1-(1-p_{\text{iter}})^{\tau}
	\;\ge\;
	\Omega\!\bigl(\min\{1,\;\tau\,n^{-(2k+2)}\}\bigr).
	\]
	
	\noindent
	%The hidden constant depends only on the fixed block size \(k\).
\end{proof}



\begin{lemma} \label{lem:DPO-epoch}
	\textcolor{blue}{Let an ageing epoch reach the junction $J$.}% (which occurs with high probability
	%by Lemma~\ref{lem:junction-first-density-Bk}; we condition on this event
	%throughout).}
	Then, with probability $\new{p_D}=\Omega(\min\{1,\,\tau\,n^{-\beta}\})$, the epoch performs a jump placing the individual onto the deep-prefix plateau.
\end{lemma}

At $J$, the \textcolor{black}{ones-density} of the suffix already exceeds $\rho$ with high probability.
Thus it suffices to flip the prefix to $0^L *$.
Under the same high-Hamming-distance event ensuring $M(x)\ge cn/20$, the probability that the first $L$ flips hit positions $\{1,\dots,L\}$ is at least $(c/20)^L = n^{-\beta}$.
Evaluating over $\tau$ iterations gives the lemma.

\begin{proof}[Proof of Lemma \ref{lem:DPO-epoch}]
	At the beginning of the epoch the reference is the best string from the previous
	epoch (from $J$ or $\mathcal{B}_k$), and the new parent $x$ itself lies in~$J$.
	Hence,
	\[
	|x|_1\;=\;0.9n-1 \;+\; |x_{\text{suffix}}|_1.
	\]
	The suffix has length $0.1n+1$ and is uniform; its number of ones is
	$\mathrm{Bin}(0.1n+1,\tfrac12)$ with expectation $0.05n+O(1)$.
	By a Chernoff bound
	\(
	|x_{\text{suffix}}|_1\ge 0.04n
	\)
	with probability $1-e^{-\Omega(n)}$.
	Thus, already before any mutation,
	\[
	|x|_1
	\;\ge\;
	0.9n-1 + 0.04n
	\;>\;
	0.94n
	\;>\;
	\rho n
	\qquad(\rho<0.9).
	\tag{1}
	\]
	
	%\textbf{Step 1 – Linear mutation potential.}
	As in the block-jump argument we have
	\(
	\Pr\!\bigl(\mathrm{HD}(x,best)\ge n/20\bigr)\ge\tfrac12
	\)
	(owing to the uniform suffix). We condition on this event, so we have
	 \(M\ge cn/20\) throughout the epoch.
	
	%\textbf{Step 2 – Per-iteration flip of the \(L\) prefix bits.}
	%Let \(L=\lfloor\log_2 n\rfloor\).
	To generate a $\mathcal{D}_\rho$ optimum it is necessary to flip the 
	\(L\) leading ones from $1$ to $0$ before FCM stops.
	Sampling without replacement,
	\begin{align*}
		p_{\text{prefix}}
		%&=\Pr\bigl(\{1,\dots,L\}\subseteq\text{ first }L\text{ flips}\bigr)
		&=\frac{(M \cdot (M-1)\cdot  ... \cdot (M-L))}{(n \cdot (n-1)\cdot  ... \cdot (n-L))}
		\;\ge\;
		(M/n)^{L}
		\;\\&\ge\;
		(cn/20)^{L}
		=n^{-\beta}.
	\end{align*}
	where $\beta:=\log_{2}\frac{20}{c}$.
	
	%\textbf{Step 3 – Density constraint is automatic.}
	\textcolor{blue}{When those $L$ flips occur,
	since the $L$-bit prefix of the new parent is a uniformly random string
	(from the reinitialisation), it contains approximately $L/2$ one-bits in
	expectation; the total number of ones in the string therefore drops by at
	most $L/2 \le \tfrac12\log_2 n$ when these are flipped.
	Thus using (1) and $L = o(n)$ we get,}
	%\[
	$|x|_1-L/2\;\ge\;0.94n-\tfrac12\log_2 n
	\;\ge\;\rho n$
	%\quad\text{
	for all sufficiently large $n$.
	%\]
	Hence the density requirement \(|x|_1\ge\rho n\) is already satisfied
	once the prefix flips have happened. Thus no additional tail flips are needed, and
	FCM immediately accepts a $\mathcal{D}_\rho$ search point.
	
	%\textbf{Step 4 – Per-epoch probability.}
	Since each iteration satisfies \(p_{\text{iter}}\ge n^{-\beta}\).
	Across the \(\tau\ge n\) iterations of the epoch,
	\[
	p_D
	\;\ge\;
	1-(1-p_{\text{iter}})^{\tau}
	\;\ge\;
	%\min\!\bigl\{1,\;\tau\,p_{\text{iter}}\bigr\}
	1-(1-n^{-\beta})^\tau 
	\;=\;
	\Omega\!\bigl(\min\{1,\;\tau\,n^{-\beta}\}\bigr),
	\]
	completing the proof.
\end{proof}


\textcolor{blue}{Each \emph{ageing epoch} begins with a uniformly random reinitialisation and
runs for $\tau$ iterations; by Lemma~\ref{lem:junction-first-density-Bk} it
ascends to the junction $J$ with high probability before reaching any
off-slope structure.  We henceforth condition on epochs reaching $J$.}
Each ageing epoch evaluates $\Theta(n\tau)$ individuals.
By Lemma~\ref{lem:junction-first-density-Bk}, epochs reach $J$ with high probability, and by \textcolor{black}{lemmata~\ref{lem:BJO-epoch} and~\ref{lem:DPO-epoch}}, an epoch discovers each off-slope structure with probabilities $p_B$ and $p_D$, respectively.
\textcolor{blue}{Since $B$ (entering $\mathcal{B}_k$) and $D$ (entering $\mathcal{D}_\rho$) are
mutually exclusive within a single epoch, the algorithm must discover them in
\emph{different} epochs.
By a coupon-collector argument (we need to ``collect'' both coupons $B$ and
$D$, which arrive independently across epochs with probabilities $p_B$ and
$p_D$ respectively), the expected number of epochs until both optima have been seen
is $O(1/p_B + 1/p_D)$.}
Multiplying by the cost per epoch yields the bound stated in Theorem~\ref{thm:main-tau-final}.

\begin{proof}[Proof [of Theorem~\ref{thm:main-tau-final}]]
	%\textbf{1.  Evaluation cost of one epoch.}
	%Each iteration evaluates at most \(M(x)\le c\,n\;(=n/3)\) offspring, so an
	%epoch of length \(\tau\) uses
	%\[
	%\text{eval}_{\text{epoch}}\le c\,n\,\tau=\Theta(n\,\tau).
	%\]
	
	%\textbf{2.  Per–epoch success probabilities.}
	By Lemma~\ref{lem:junction-first-density-Bk} we reach the junction $J$ with probability \ \(1-o(1)\).
	Starting an epoch in the junction, %\(J\) (prob.\ \(1-o(1)\)), 
	the two lemmata give
	%\[
	$p_B=\Omega\!\bigl(\min\{1,\tau\,n^{-(2k+2)}\}\bigr)$ and
	$p_D=\Omega\!\bigl(\min\{1,\tau\,n^{-\beta}\}\bigr)$,
	%\]
	and \(B\) and \(D\) are mutually exclusive within one epoch. Thus the expected number of epochs to sample from both sets of optima is:
	
	%\textbf{3.  Expected number of epochs to see both optima.}
	%A coupon–collector bound yields
	\[
	%\mathbb{E}[T_{\mathrm{both}}]
	%\;\le\;
	%\frac1{p_B}+\frac1{p_D}
	%\;=\;
	O\!\Bigl(
	\max\{1,n^{2k+2}/\tau\}
	+\max\{1,n^{\beta}/\tau\}
	\Bigr).
	\]
	
	%\textbf{4.  Total expected evaluations.}
	Multiplying by the cost of each epoch we get,
	\begin{align*}
		\mathbb{E}[\text{evals}]
		&\le
		\Theta(n\,\tau)\;
		O\!\Bigl(
		\max\{1,n^{2k+2}/\tau\}
		+\max\{1,n^{\beta}/\tau\}
		\Bigr)
		\\&=
		O\!\bigl(
		n\,\tau
		+n^{2k+3}
		+n^{\beta+1}
		\bigr).
	\end{align*}
\end{proof}

\textcolor{black}{The static hypermutation would yield a same upper bound with a similar analysis except for the value of $\beta$ which would be smaller and for the range of $c$ which would be more restricted. Both modifications are necessary due to the larger mutation size of the static hypermutation compared to IPH for a given $c$ value. The \lastBest~IPH cannot make the jumps from the junction since any current junction solution would be the last best seen, yielding a mutation potential of one $1$. $\mathcal{B}_k$ and $\mathcal{D}_\rho$ sets cannot be accessed from the junction via mutation but can still be sampled at any reinitialization with $2^{-\Omega(n)}$ probability.}
\textcolor{purple}{We conjecture that power-law mutation operators such as those employed by the Fast-IA~\cite{Fast-AIS} and the Fast-EA~\cite{Fast-GA} would also succeed on $\mathrm{LTJ}_{k,\rho}$ if equipped with and ageing operator such as the algorithms introduced in~\cite{Fast-TEVC}, and the same should hold for the Rank-GA~\cite{RankGA}. However, we believe that the Rank-GA with ageing would fail on \textsc{Cliff}$_d$ because of its high mutation rate at the bottom of the second slope. On the other hand, the Fast-IA~\cite{Fast-AIS} and the Fast-EA~\cite{Fast-GA}, which have been shown to be parameterisable such that they behave equivalently, can quite efficiently (a $\sqrt{n}$ factor slower than the \firstBest IPH i.e., in expected time $O(n^{5/2} \log n$))  optimise \textsc{Cliff}$_d$ functions with gaps of linear size~\cite{Fast-TEVC}. However, they have not been shown to be efficient for \textsc{Cliff}$_d$  functions with sublinear gap sizes. 
Stagnation detection (?!?) algorithms~\cite{} may also be efficient for $\mathrm{LTJ}_{k,\rho}$ if equipped with the ageing operator. It is not obvious, though, how the operator would be incorporated in such an algorithm. Thus, to date the \firstBest IPH is the best known algorithm to optimise both Cliff-like functions while at the same time optimise functions where large mutations are necessary to escape from local optima. An open question is whether it may be possible to avoid the use of ageing altogether for $\mathrm{LTJ}_{k,\rho}$ by using the recently introduced population-based EAs with low selective pressure~\cite{TELO,GECCO}. Low selective pressure may avoid the entire population converging on one of the two local optima shortly after it is identified.} 
